\section{The “Prophecies”  of  Rene Guenon}

In the conclusion of \emph{Introduction to the Study of Hindu Doctrines}, originally published in 1921, \textbf{Rene Guenon} offers some predictions, which he emphasises are not to be considered ``prophecies'', for the spiritual direction of the West. He points out that the obstacle preventing a re-orientation toward Tradition is that most Westerners regard the decadence that so disgusts Guenon as ``progress'', so that there exists little desire for such a change.

He holds out the chance that ``people will one day begin to notice that things which now appear all-important are unable to yield the results expected of them''. Unfortunately, he says that this ``disillusion'' will be of a merely a sentimental character, hence the intellectual source for genuine change will have to come from the outside. What the West regards as progress -- material and scientific advancement---is accompanied by ``retrogression in the purely speculative and intellectual order''. (This theme is taken up by Evola in \emph{Pagan Imperialism}.)

Guenon then describes three possible scenarios, though he hedges on the timeframe: from ``soon'' to ``several centuries'' from now.

\paragraph{Degeneration}


The worst scenario would occur if nothing were introduced to change course. Then the West, ``abandoned to its own fate, would sink into the lowest forms of barbarism''. This can only sound implausible to those who believe in continual progress, but Guenon points out many other civilisations have declined and even disappeared.

\paragraph{Assimilation}


The second possibility is that outsiders from the East would rescue the West through assimilation. And this scenario has two options: assimilation by consent or by force. Guenon adds: ``assuming that the thing were possible and that the East were willing to do this.'' Of course, in 1945, Guenon could not have anticipated the sudden and massive influx of Muslims into Europe within two generations later, so it is certainly possible to consider that as fulfilling the conditions for assimilation.

Guenon regards this scenario as preferable to the first and warns that Western ``prejudice'' (read: ``racism'') would blind them to prevent recognition of this. The next passage is worth quoting in full:

\begin{quotex}
under such circumstances there would doubtless be a transition period of extremely painful ethnical revolutions, which are difficult to picture but which in their final result would be of a nature to compensate for the damage certain to be sustained during a catastrophe of this kind; but in that case the West would have had to forego its own character and would find itself absorbed purely and simply.
\end{quotex}


This judgment may be premature, but this has the ring of ``prophecy''.

\paragraph{Transformation}


The Transformation scenario is regarded by Guenon as more favourable from the Western point of view, though it still would involve the disappearance of the West as currently constituted. This is the ``return to true and normal intellectuality''. However, unlike the assimilation scenario, which involves the imposition of the transformation under duress from outsiders, this scenario supposes the West brings it about voluntarily and spontaneously.

It's amazing how many currently burning issues were anticipated by Guenon. For example, the project of the West is to export enlightenment values --- i.e., science, democracy, liberalism --- to the East in the hope to alter the mentality of Easterners. Of this, Guenon writes: ``Only a delusion and a blindness begotten of the most ridiculous prejudice could allow a man to believe that the Western mentality can win over the East.'' Western ``intellectuals'' only come into contact with Easterners who are Westernized by the education system, so they get a skewed view of things. Yet, as events of the past few years show, this Westernization is often superficial even among the most educated Easterners. This seems to cause nothing but puzzlement in Western minds.

It might seem odd that the concluding chapter on a book about Hindu doctrines is all about the ``West''. Obviously, the intent is not for the West to become Hindu, but rather for the West to use the traditional doctrines as revealed in the Vedanta in order to become what it truly is.

Guenon ends the chapter with a deeper review of the three alternatives he mentioned. ``Everything clearly depends on the mental state of the Western world at the moment when it reaches the furthest term of its present civilization.'' If the mental state remains as it currently is -- viz., based on non-traditional values which both Evola and Guenon have analysed in various books -- then the first alternative will prevail. This is because there would be nothing left in the Western mentality to accept Tradition, either from within, but also by way of assimilation.

The second alternative -- assimilation -- would still require an intellectual kernel, even if confined to a numerically small elect. This elect would have to be strong enough to serve as intermediaries to guide the West back to the sources of ``true intellectuality''. This does not require the conscious awareness of the masses to be effective. Guenon warns that this alternative ``would not be free from certain unpleasant features''. These unpleasant features are a matter for speculation, but probably not difficult to discern.

Now, if the elect were strong enough and had sufficient time to operate on its own, then the third possibility could be realised. It would have to act as a ``leaven'', slowly influencing the Western world. This does not require the consent of the mass of Westerners, nor even for the generality of so-called intellectuals. ``It is enough, therefore, as a start, for a few individuals to understand the need for such a change, but of course on condition that they understand it truly and thoroughly.''

Guenon points out: ``A return to a traditional civilization, both in principle and in respect of the whole body of institutions, is evidently the base condition for the transformation we have been speaking about.'' Clearly, a return to tradition is the most essential purpose of the elect, but how to bring this about. A concrete example would be helpful.

As an example, Guenon provides: ``We can only say that the Middle Ages afford us an example of a traditional development that was truly Western; ultimately it would be a case not purely and simply of copying or reconstructing what existed then, but of drawing inspiration from it in order to bring about an adaptation to suit the actual circumstances. \emph{If there exists a ``Western tradition'', that is where is must be looked for, and not in the fantasies of occultists and pseudo-esotericists.}'' My emphasis. Lest anyone thing Guenon is being idiosyncratic here, compare to Evola: ``the grandeur of Rome, having risen from the forces of the Nordic Aryans, created the last, great, universal period in the West, the feudal-imperial civilization of the Middle Ages.'' (from ``Pagan Imperialism'') And again, ``The power of a new Middle Ages is needed.'' Sir John Woodroffe also regards that era as traditional in comparison to the modern West. (``The World as Power'')

Unless there is a secret ``elect'' at work, why is it that no one looks to that model? Instead, we have the ``Nouvelle Droite'', Asatru and similar neo-pagan reconstructions, revivals of Germanic occultist organisation, and so on. The Middle Ages seem to be an embarrassment to everyone. To the contrary, to the extent that one cannot recognize the traditional elements at play there, to an even greater extent, it is unlikely that traditional elements can be recognized anywhere at all, never mind serve as a leaven for transformation.

\flrightit{Posted on 2008-01-10 by Cologero}

\begin{center}* * *\end{center}

\begin{footnotesize}
\begin{sffamily}
\texttt{Saladin on 2013-06-03 at 20:26 said:}

As an Easterner and a follower of Guenon I have to say that his idea of the ``East'' was a bit idealized if not purely romanticized. The fact is that the process of degeneration has effected the East as much as the West. ``The Desert approaches! Woe to the one whose desert is within!''

\hfill

\texttt{August on 2013-06-04 at 05:40 said:}

Saladin, looking around today and taking first impressions, we could reasonably venture that Guenon exaggerated the fidelity of the East. Yet he did say, with regard to the extensive Westernisation of the East already in motion during his lifetime, that the only Easterners that truly matter are those that preserve their intellectuality, and that their numbers needn't be very high to ensure the superiority of the East.

What is the likelihood that these small groups, `elites' perhaps, persist untainted in the East? Probably very high. Consider the following:

-Despite the crises wracking the Islamic world at the moment, religion persists as a centrepiece of life for the average Muslim. The wholesale repudiation (or perversion) of one's faith is, despite everything, a relatively rare phenomenon in Dar al-Islam, while being the dominant characteristic of the West. Even the most grotesque `fundamentalist' reactions to the threatening spectre of modernity are, in the end, spiritually aspirational, even if stunted and deformed. Why should it be so difficult to displace Islam as the nucleus of a large portion of humanity, even when it apparently costs them materially, and oppositional tendencies have already taken root in their lands?

-The Chinese, like other peoples who preserve the `sense of eternity', seem capable of reorganising their society in harmony with cyclical transitions and resulting corporeal pressures, without disintegrating or losing anything essential. What could explain this ability? That tens of millions of their fellow countrymen should die in the process hardly seems to bother them. What is death in the shadow of the eternal?

That said, Evola was right to remind us that the processes which engendered the modern West did not happen overnight, or without great turmoil. How well the East holds up in the near future remains to be seen.

\hfill

\texttt{Saladin on 2013-06-04 at 22:31 said:}

August, it's true that the process of degeneration did not happen overnight in the West (it took more than two centuries at least) but in the East it took mere decades. I agree that there are small groups of Elites that persist in the East but I hope that holds true for the West as well!

That said, the general prospects looks grim globally and I have no illusions that it will get any better.

\hfill

\texttt{Cologero on 2013-06-04 at 23:28 said:}

Saladin, are you familiar with Javadi-e Amoli?

\hfill

\texttt{Saladin on 2013-06-05 at 07:37 said:}

Cologero, having been raised in Iran I am vaguely familiar with Javadi-e Amoli. Why didyou ask?

\hfill

\texttt{Max on 2013-09-29 at 09:16 said:}

Here are some of my thoughts on our direction.

Is it destiny that all countries must become ``the west'' before the great transition, the ``reversal of the poles''? How could the east rescue the west when they seem incapable or unwilling to stop modernity even in their own countries. What must be done by a spiritual elite is the right preparations. If the whole world passes through the point to the other side at the same time, even ``globalism'' seems to have a purpose. The east can in that view not exist anymore as a physical location, but only in our hearts. It is a question of a battle between two conceptions of time, which nonetheless will both reach the same conclusion in a strange unification of opposites at the end of time. It is a passage through the still center. And what matters when that happens is what we brought with us, the result of everything, then distilled into a new seed.

\hfill

\texttt{Cobbler on 2015-07-30 at 13:13 said:}

Hello,

@Saladin, you should read ``West'' as the representation of the anti-traditional world, the modern mind, and the ``East'' as the representation of the Traditional Order ; not strictly ``geographically speaking beyond the locations''.

Guénon wrote it, this ``apparent chaos'' will eventually spread all over the world ; looking around us, we can't deny it, even if we have to be careful how you ``interpret'' those changing effects ;

indeed, as appointed and mentioned by @Cologero about the Islamic world and also the Chinese world which indeed have more similarities than it seems but this is another topic that will bring us too far afield.

Best Regards.

\hfill

\texttt{Cobbler on 2015-07-30 at 13:22 said:}

erratum: mentioned by @August

\hfill

\texttt{August on 2015-07-31 at 07:37 said:}

Hi there Cobbler,

I'm aware that we can read the two terms in that sense, but in this case I was using them with their literal geographic or ethnic meaning since I was discussing the present state of the world and her people.

Needless to say the state of most easterners one meets today is quite apparent; like most others they seem to me rather ordinary people and not much else, religious or not. Let there be no illusions.

\hfill

\texttt{Harun on 2016-01-29 at 19:39 said:}

Sadly I think the first predictions is the direction the west is following

The current muslim immigrants do not belong to the traditional and highest form of Islam, Whabites are anti-traditional extoerical fanatics who despise sufism.

Truth is, Islam is being subverted too.

\hfill

\end{sffamily}
\end{footnotesize}

